\documentclass[11pt]{article}
\usepackage{amsmath, amssymb}
\usepackage{geometry}
\usepackage{enumitem}
\usepackage{array}
\usepackage{booktabs}
\geometry{margin=0.7in}

\title{Ryan Algebra II Final Practice Grading}
\author{}
\date{}

\begin{document}
\maketitle

\section*{Grading Note}
The submitted reference image matches \textit{Algebra II Final Targeted Practice}, not \textit{Algebra II Final Topics Review}. I graded the visible work against the targeted practice sheet. Problems 5--12 were not visible in the provided reference image, so they are marked blank.

\section*{Score Summary}
\[
\text{Visible work score: } 3.9/4
\qquad
\text{Overall submitted score: } 3.9/12
\]

\begin{center}
\begin{tabular}{c c l}
\toprule
Problem & Credit & Note\\
\midrule
1 & $1/1$ & Exact inverse-trig values are correct.\\
2 & $1/1$ & Correct factorization and solution set.\\
3 & $1/1$ & Correct cosine values and angles.\\
4 & $0.9/1$ & Features are essentially correct; amplitude should be positive $3$.\\
5--12 & $0/8$ & No submitted work visible in the reference image.\\
\bottomrule
\end{tabular}
\end{center}

\section*{Recreated Work and Corrections}
\begin{enumerate}[leftmargin=*, itemsep=1.1em]
  \item Ryan's visible work:
  \[
  \tan\theta=-\frac{5}{12},\quad 5^2+12^2=169,\quad \sqrt{169}=13,
  \quad \sec\theta=\frac{13}{12}
  \]
  \[
  \sin\!\left(\cos^{-1}\!\left(-\frac{7}{25}\right)\right)=\frac{24}{25},
  \qquad
  \cot\!\left(\sin^{-1}\!\left(\frac45\right)\right)=\frac34
  \]
  Correct answer:
  \[
  \boxed{\frac{13}{12}},\qquad \boxed{\frac{24}{25}},\qquad \boxed{\frac34}
  \]
  Credit: full.

  \item Ryan's visible work:
  \[
  2\cos^2\theta-\cos\theta-1=0,\qquad
  (2\cos\theta+1)(\cos\theta-1)=0
  \]
  \[
  \cos\theta=-\frac12 \quad\text{or}\quad \cos\theta=1
  \]
  Correct answer:
  \[
  \boxed{\theta=0,\ \frac{2\pi}{3},\ \frac{4\pi}{3}}
  \]
  Credit: full. Since $0\le \theta<2\pi$, $2\pi$ is not included.

  \item Ryan's visible work:
  \[
  \sqrt3\sec\theta+2=0,\qquad
  \sqrt3\sec\theta=-2,\qquad
  \sec\theta=-\frac{2}{\sqrt3}
  \]
  \[
  \cos\theta=-\frac{\sqrt3}{2}
  \]
  Correct answer:
  \[
  \boxed{\theta=\frac{5\pi}{6},\ \frac{7\pi}{6}}
  \]
  Credit: full.

  \item Ryan's visible work:
  \[
  y=-3\cos\left(2x+\frac{\pi}{2}\right)-1
  =
  -3\cos\left(2\left(x+\frac{\pi}{4}\right)\right)-1
  \]
  Recreated features:
  \[
  \text{amplitude written as } -3,\qquad
  \text{period }=\pi,\qquad
  \text{phase shift }=-\frac{\pi}{4},
  \]
  \[
  \text{midline }y=-1,\qquad
  \text{range }[-4,2]
  \]
  Correct answer:
  \[
  \boxed{\text{amplitude }=3},\qquad
  \boxed{\text{period }=\pi},\qquad
  \boxed{\text{phase shift }=-\frac{\pi}{4}}
  \]
  \[
  \boxed{\text{midline }y=-1},\qquad
  \boxed{\text{range }[-4,2]},\qquad
  \boxed{\text{one full period: }\left[-\frac{\pi}{4},\frac{3\pi}{4}\right]}
  \]
  Credit: nearly full. The vertical coefficient is $-3$, but amplitude is the positive distance $3$.

  \item No visible submitted work.
  \[
  \boxed{y=2\sin\left(\frac12\left(x-\frac{\pi}{2}\right)\right)-3}
  \]

  \item No visible submitted work.
  \[
  c=\sqrt{16^2+21^2-2(16)(21)\cos58^\circ}\approx 18.5
  \]
  \[
  K=\frac12(16)(21)\sin58^\circ\approx 142.5
  \]
  \[
  \angle B\approx 75^\circ
  \]
  The larger angle is $\boxed{\angle B}$.

  \item No visible submitted work. Since
  \[
  \sin Q=\frac{24\sin42^\circ}{18}\approx 0.892,
  \]
  two triangles are possible:
  \[
  \boxed{Q\approx 63^\circ,\ R\approx 75^\circ,\ r\approx 26.0}
  \]
  \[
  \boxed{Q\approx 117^\circ,\ R\approx 21^\circ,\ r\approx 9.7}
  \]

  \item No visible submitted work. The bearings differ by $90^\circ$, so
  \[
  d=\sqrt{5.2^2+7.1^2}\approx \boxed{8.8\text{ miles}}.
  \]

  \item No visible submitted work.
  \[
  x^2-10x+8y-7=0
  \]
  \[
  (x-5)^2=-8(y-4)
  \]
  This is a parabola with
  \[
  \boxed{\text{vertex }(5,4)},\quad
  \boxed{\text{focus }(5,2)},\quad
  \boxed{\text{directrix }y=6},\quad
  \boxed{\text{axis }x=5}.
  \]

  \item No visible submitted work.
  \[
  9x^2+4y^2-54x+32y+109=0
  \]
  \[
  \frac{(x-3)^2}{4}+\frac{(y+4)^2}{9}=1
  \]
  This is an ellipse with
  \[
  \boxed{\text{center }(3,-4)},\qquad
  \boxed{\text{major axis length }6},\qquad
  \boxed{\text{minor axis length }4}.
  \]

  \item No visible submitted work.
  \[
  25x^2-16y^2-200x-32y-16=0
  \]
  \[
  \frac{(x-4)^2}{16}-\frac{(y+1)^2}{25}=1
  \]
  This is a hyperbola with
  \[
  \boxed{\text{center }(4,-1)},\quad
  \boxed{\text{vertices }(0,-1),(8,-1)}
  \]
  \[
  \boxed{\text{foci }(4-\sqrt{41},-1),(4+\sqrt{41},-1)},\quad
  \boxed{y+1=\pm\frac54(x-4)}.
  \]

  \item No visible submitted work.
  \[
  \boxed{\frac{(y-3)^2}{25}-\frac{(x+2)^2}{16}=1}
  \]
  Since $c^2=a^2+b^2=25+16=41$,
  \[
  \boxed{\text{foci }(-2,3-\sqrt{41}),\ (-2,3+\sqrt{41})}.
  \]
\end{enumerate}

\end{document}
